\section{Struttura e configurazione di una
Challenge}\label{struttura-e-configurazione-di-una-challenge}

In BlueAgent, la definizione di una challenge adotta l'approccio
\emph{Infrastructure-as-Code} (IaC): l'intera topologia viene descritta
mediante uno o più file dichiarativi in formato YAML.

Questo file specifica quali container compongono l'ambiente, le risorse
hardware assegnate a ciascuno di essi ed i comandi di readiness e di
verifica necessari per convalidarne la patch.

La configurazione della challenge avviene in due step:

\begin{itemize}
\tightlist
\item
  Specifica delle immagini Distrobuilder tramite file di configurazione.
\item
  Configurazione della challenge vera e propria. Vengono specificati qui
  i container da istanziare e con quali immagini.
\end{itemize}

\subsection{Struttura del File di Configurazione della
Challenge}\label{struttura-del-file-di-configurazione-della-challenge}

La struttura della specifica di una challenge è articolata nelle
seguenti sezioni principali:

\begin{enumerate}
\def\labelenumi{\arabic{enumi}.}
\item
  \textbf{Metadati e Identificazione:}

  \begin{itemize}
  \tightlist
  \item
    \passthrough{\lstinline!challenge\_id!}: Identificatore alfanumerico
    univoco globale della sfida (ad esempio
    \passthrough{\lstinline!patch-web-sqli!}).
  \item
    \passthrough{\lstinline!team\_k!}: Indice numerico deterministico
    del team (\(1 \le k \le 254\)).
  \end{itemize}
\item
  \textbf{Topologia dei Container
  (\passthrough{\lstinline!containers!}):} Rappresenta la lista delle
  macchine che compongono la sfida. Per ciascun container vengono
  dichiarati:

  \begin{itemize}
  \tightlist
  \item
    \passthrough{\lstinline!name!}: Nome logico della macchina
    all'interno della challenge (ad esempio
    \passthrough{\lstinline!web!} o \passthrough{\lstinline!database!}).
  \item
    \passthrough{\lstinline!image!}: Alias dell'immagine Incus da
    istanziare.
  \item
    \passthrough{\lstinline!ip\_suffix!}: Ultimo ottetto dell'indirizzo
    IP statico assegnato al container
    (\passthrough{\lstinline!10.k.1.X!}), dove
    \passthrough{\lstinline!10.k.1.1!} è riservato al gateway.
  \item
    \passthrough{\lstinline!limits\_cpu!} e
    \passthrough{\lstinline!limits\_memory!}: Vincoli hardware imposti
    tramite cgroup.
  \end{itemize}
\item
  \textbf{Runner di Verifica Effimero
  (\passthrough{\lstinline!verifier!}, opzionale):} Definisce un
  container di test dedicato, istanziato esclusivamente durante la fase
  di verifica all'interno della rete OVN (ad esempio con suffisso IP
  \passthrough{\lstinline!250!}).
\item
  \textbf{Comandi di Readiness
  (\passthrough{\lstinline!readiness\_commands!}):} Serie di controlli
  preliminari eseguiti all'avvio della sandbox per accertarsi che i
  servizi interni abbiano completato la fase di bootstrap, prima di
  invocare i test veri e propri. Ciascun comando include intervalli di
  retry e policy di timeout.
\item
  \textbf{Comandi di Verifica
  (\passthrough{\lstinline!verify\_commands!}):} Rappresentano la suite
  di convalida che certifica la patch. BlueAgent supporta due modalità
  operative distinte:

  \begin{itemize}
  \tightlist
  \item
    \passthrough{\lstinline!target: container!}: verifica in White Box
    eseguita direttamente all'interno dell'istanza bersaglio (ad esempio
    per verificare lo stato di un servizio con
    \passthrough{\lstinline!systemctl!}).
  \item
    \passthrough{\lstinline!target: verifier!}: verifica in Black Box
    eseguita dal container tester per simulare l'attacco di rete o
    verificare l'integrità del servizio attraverso la LAN virtuale.
  \end{itemize}
\end{enumerate}

\textbf{Esempio di Configurazione YAML:} Di seguito è riportato un
esempio sintetico di descrittore dichiarativo per una challenge
multi-tier composta da un'applicazione web e un database:

\begin{lstlisting}
challenge_id: "patch-web-sqli"
verify_internet_access: true

containers:
  - name: "database"
    image: "debian-12-mysql"
    ip_suffix: 3
    limits_cpu: "1"
    limits_memory: "512MB"

  - name: "web"
    ...

verifier:
  image: "verifier-python"
  ...

readiness_commands:
  - target: "container"
    container_name: "database"
    cmd: ["mysqladmin", "ping", "-h", "127.0.0.1"]
    retry: 2
    timeout: 15

verify_commands:
  - target: "container"
    container_name: "web"
    cmd: ["systemctl", "is-active", "--quiet", "flask-app"]
    timeout: 5
  - target: "verifier"
    ...
\end{lstlisting}

\textbf{Parsing e Risoluzione Dinamica degli Indirizzi IP nei Comandi:}
Poiché la specifica YAML viene scritta dall'autore della sfida in modo
indipendente dallo specifico team, i comandi non possono contenere
indirizzi IPv4 cablati staticamente. BlueAgent implementa un meccanismo
di template interno:

\begin{itemize}
\tightlist
\item
  \passthrough{\lstinline!\{\{ip:<name>\}\}!}: viene tradotto
  nell'indirizzo IPv4 completo
  \passthrough{\lstinline!10.<team\_k>.1.<ip\_suffix>!} associato al
  container con nome logico \passthrough{\lstinline!<name>!}.
\item
  \passthrough{\lstinline!\{\{subnet\}\}!}: viene sostituito con il
  prefisso di rete assegnato al team
  (\passthrough{\lstinline!10.<team\_k>.1!}).
\item
  \passthrough{\lstinline!\{\{verifier\_ip\}\}!}: corrisponde all'IP
  assegnato al runner di test
  (\passthrough{\lstinline!10.<team\_k>.1.<ip\_suffix>!}).
\item
  \passthrough{\lstinline!\{\{team\_k\}\}!}: viene sostituito
  dall'identificatore numerico del team.
\end{itemize}

\begin{center}\rule{0.5\linewidth}{0.5pt}\end{center}

\subsection{Configurazione delle Immagini con
Distrobuilder}\label{configurazione-delle-immagini-con-distrobuilder}

\begin{quote}
\passthrough{\lstinline!Distrobuilder!}: È lo strumento open source
ufficiale per creare immagini di sistema per Incus ed LXC a partire da
alberi di pacchetti e distribuzioni upstream ufficiali.
\end{quote}

La configurazione di un'immagine in Distrobuilder avviene mediante un
file dichiarativo. I campi principali sono i seguenti:

\begin{enumerate}
\def\labelenumi{\arabic{enumi}.}
\item
  \textbf{Sorgente ed Immagine (\passthrough{\lstinline!image!} e
  \passthrough{\lstinline!source!}):} Indica quale immagine usare come
  base.
\item
  \textbf{Gestione dei Pacchetti (\passthrough{\lstinline!packages!}):}
  Definisce l'elenco dei software da installare all'interno del rootfs
  per preparare l'ambiente. È possibile anche rimuovere pacchetti
  preinstallati nell'immagine.
\item
  \textbf{Fasi di Hook e Predisposizione Rete
  (\passthrough{\lstinline!actions!}):} Consente di eseguire script
  durante le fasi intermedie di generazione del filesystem. È essenziale
  per predisporre le directory dei demoni di rete o specificare
  configurazioni interne al sistema.
\end{enumerate}

All'interno dell'architettura di BlueAgent, il sottosistema
\passthrough{\lstinline!ImageService!} invoca Distrobuilder per generare
il tarball da importare nello storage pool ZFS.
